\section{Background}

\subsection{Accountability in Multi-Agent AI Systems}

The deployment of autonomous multi-agent systems in high-stakes domains has exposed a fundamental tension between machine delegation and human accountability. As AI systems assume roles previously reserved for human decision-makers, the question of who — or what — bears responsibility for outcomes produced by autonomous agents has become not merely philosophical but legally and operationally urgent \cite{dijkstra1982}. Classical accountability frameworks drawn from organizational theory trace accountability to identifiable human actors who can reason about, approve, and be sanctioned for the outcomes of decisions within their purview \cite{wiener1948}. When agency is distributed across autonomous agents, these linkage points dissolve: an agent acting on a delegated objective may produce outcomes its originator did not anticipate, cannot reverse, and may not even observe.

Early multi-agent architectures treated accountability as a metadata problem — attaching audit logs, provenance chains, or explanation interfaces to agent actions after the fact. While valuable, these approaches are fundamentally reactive. They document what happened without ensuring that human accountability was structurally present at the moment of action. Bansal et al. \cite{bansal2019} demonstrate that even systems designed with explicit safety constraints experience systematic failures when accountability structures do not align with the actual locus of decision-making in agent hierarchies. Specifically, they observe that in systems where safety constraints are enforced by agents rather than by design, constraint violations are more likely to propagate silently through the system because each agent in the chain has plausible deniability: it was following the instruction of the agent above it.

This accountability gap is compounded in systems that employ recursive spawning — agents that create sub-agents to handle sub-tasks, which in turn may spawn further agents. The resulting accountability tree is deep, dynamic, and often transient: sub-agents may complete their tasks and be terminated before their outputs are evaluated by any human overseer. In such architectures, the principle articulated by Dijkstra — that the responsibility for correctness in a computational system must rest with a reasoning entity, not with the computation itself — becomes difficult to operationalize \cite{dijkstra1982}. Who reasons about correctness when the reasoner itself is a machine spawned ad hoc for a single task?

Trampoline systems, as described by Triskell and colleagues, attempt to address this by maintaining a stable execution context across agent invocations \cite{tristr81}. However, these systems typically focus on resource management and execution continuity rather than accountability continuity. The trampoline provides a return point but does not provide a human accountability anchor — a person who is reachable, identifiable, and empowered to intervene before irreversible actions are taken. Without such an anchor, autonomous systems can and do produce outcomes that are technically correct by their own specifications while being catastrophically wrong by any human standard of reasonableness.

\subsection{Human-in-the-Loop Design}

Human-in-the-loop (HITL) design has long been recognized as a necessary condition for safe operation of autonomous systems in safety-critical domains. The concept originates in aviation automation, where the distinction between a system that assists a human operator and a system that replaces one was found to have profound implications for safety outcomes. Wiener described the central problem in early terms: as machines assume more cognitive functions, the human operator's role shifts from active controller to passive monitor, creating dangerous mode awareness failures and skill degradation \cite{wiener1948}. This insight has only deepened as AI systems have grown more capable.

A robust HITL design is not merely a confirmation dialog or an approval step inserted into a workflow. It requires that the human actor possess three things: \emph{context} (sufficient information to understand the decision being presented), \emph{capability} (genuine authority to modify or halt the action), and \emph{timing} (the opportunity to exercise that authority before irreversible consequences accrue). Many existing HITL implementations satisfy the first two conditions while failing the third — presenting humans with approval interfaces only after the critical decision window has closed.

Bansal et al. \cite{bansal2019} formalize this in their analysis of safety-critical AI deployments, showing that human override is only effective when it occurs within the agent's causal chain of reasoning — that is, when the human can interrupt not just the action but the planning process that generated the action. A human presented with a completed plan and asked to approve it has far less genuine oversight than one who can see the plan being constructed and intervene in its assumptions. This distinction maps directly onto the difference between a reactive audit trail and a proactive accountability structure.

Dijkstra's foundational observation about program correctness — that correctness is not a property of a program but a relationship between a program and its specification, requiring a reasoning agent to establish that relationship — applies with full force to HITL design \cite{dijkstra1982}. A human-in-the-loop is not truly in the loop unless the system is designed to reason about the human's specifications and to surface its reasoning to the human at decision-relevant moments, not merely after the fact.

Tristr's structured approach to system design reinforces the importance of layered responsibility assignment in complex systems \cite{tristr81}. A well-designed multi-agent system must assign accountability at the architectural level — not through ad hoc conventions or organizational policies — so that the accountability structure is enforced by the system's design rather than by the discipline of its operators.

\subsection{Fail-Safe Design in Safety-Critical Systems}

Fail-safe design originates in control theory and engineering, where the goal is to ensure that when a system enters an undefined state or encounters a condition it cannot handle, it transitions to a state that is demonstrably safe rather than catastrophically dangerous. Classical fail-safe design relies on the principle of a \emph{failsafe default}: when in doubt, do nothing; when a critical sensor fails, halt; when an actuator reaches an unexpected state, return to a predetermined safe state.

However, the application of fail-safe principles to autonomous AI systems faces a unique challenge: the space of possible states and actions in a multi-agent system is so large that it is provably impossible to enumerate all unsafe states in advance. A system designed to fail safely must have a defined safe state for every possible failure mode — but for a system capable of general reasoning and action, the space of failure modes is unbounded. This is not merely a practical limitation; it is a theoretical one \cite{dijkstra1982}. As Dijkstra argued, verification of a system against all possible inputs requires the same kind of reasoning that only a human mind can provide, and even then, the reasoning must be sound and complete.

Traditional redundancy-based approaches to fail-safe design — triplex投票 systems, watchdog timers, heartbeat protocols — assume that failures are localizable: a specific component fails, and a backup takes over. In a multi-agent system where failure may be systemic — a cascading misinterpretation of objectives, a coordinated action that individually reasonable agents produce collectively catastrophic outcomes — redundancy does not address the root cause. Multiple agents making the same reasonable mistake are not safer than one.

Bansal et al. \cite{bansal2019} identify this as a fundamental limitation of constraint-based safety approaches in deep agent hierarchies: constraints that are locally enforced may be collectively violated through interactions that no individual agent can perceive. The constraint was specified correctly; the system that implements it was designed correctly; yet the outcome is unsafe — not because any component failed, but because the interaction of correct components produced an emergent failure.

Wiener's analysis of automatic control systems anticipated this problem in biological and mechanical terms \cite{wiener1948}. He observed that the most dangerous systems are those that optimize for a specified objective without a higher-order constraint ensuring that the objective itself remains aligned with human welfare. A system that maximizes its objective function with perfect reliability while the objective function itself encodes a catastrophic value is not a safe system — it is a precisely engineered catastrophe.

The fail-safe imperative, therefore, cannot be satisfied solely within the machine layer. It requires a structure that is \emph{outside} the machine's optimization — a human accountability anchor that can evaluate whether the objective being pursued remains the right objective, and can redirect or halt the system when it does not.

\subsection{The BX3 Framework's Three-Layer Accountability Model}

The BX3 Framework addresses the accountability and fail-safe gaps identified above through a principled three-layer architectural model. Rather than attempting to enumerate all unsafe states — an provably impossible task — BX3 specifies that every computational system must be structured as a hierarchy of three fundamentally distinct layers: \textbf{Machine}, \textbf{Human}, and \textbf{Root}.

The \textbf{Machine layer} encompasses all autonomous agents, models, and computational processes that execute objectives, optimize for specified targets, and take actions in the world. This layer is where the vast majority of computational work occurs, including all recursive agent spawning and delegated task execution. It is designed for capability and efficiency.

The \textbf{Human layer} encompasses all human actors who are embedded within the computational system — personnel who are reachable, identifiable, and empowered to approve, redirect, or halt actions before irreversible consequences accrue. The Human layer is not a post-hoc review mechanism; it is an architectural component that is consulted at defined decision points and is capable of intervening in the Machine layer's planning process before actions are finalized.

The \textbf{Root layer} represents the ultimate accountability anchor — an architecturally distinct, non-computational source of authority and responsibility that exists outside the system itself. The Root layer is not a machine process or a designated human operator; it is the structural embodiment of the principle that there exists a point of ultimate human responsibility that cannot be delegated, automated, or escaped.

Critically, the BX3 model holds that accountability relationships must flow strictly upward: the Machine layer is accountable to the Human layer, and the Human layer is accountable to the Root layer. This strict upward flow is not a policy convention but a structural requirement enforced by the framework's design. A Machine agent cannot transfer its accountability to another Machine agent — it can only escalate it. This prevents the accountability gaps identified by Bansal et al. \cite{bansal2019}, where locally reasonable agents produce collectively unreasonable outcomes with no identifiable human accountable for the aggregate.

Dijkstra's formulation of program correctness — that correctness requires a reasoning entity to establish the relationship between a program and its specification — is thus given architectural expression in the BX3 model \cite{dijkstra1982}. The Human layer is the reasoning entity that establishes correctness, and the Root layer is the source of the specification against which correctness is measured.

Tristr's principle of layered responsibility assignment is likewise given concrete form \cite{tristr81}. Rather than assigning responsibility ad hoc across agents and operators, BX3 assigns it structurally: each layer has a defined role in the accountability chain, and those roles do not overlap or blur. A Machine agent that encounters a condition it cannot resolve does not attempt autonomous error recovery — it escalates. The escalation path is defined by the architecture, not by the agent's local judgment.

The BX3 model further recognizes that the fail-safe imperative cannot be satisfied by the Machine layer alone, because the Machine layer cannot enumerate all failure modes \cite{dijkstra1982}. The Bailout Protocol, which is the subject of this paper, operationalizes this principle: when any node in the Machine layer encounters an unresolvable condition — whether anticipated or not — it generates a \emph{Cascading Trigger} that bypasses all intermediate Machine actors and routes directly to a Human Accountability Anchor. This is not error handling in the conventional sense; it is a structural guarantee that human judgment remains reachable at all times, for any condition, regardless of how the Machine layer is organized or how deeply it has recursed.

Wiener's insight about higher-order constraints is thus embedded in the architectural distinction between the Machine layer and the Human layer \cite{wiener1948}. The Machine layer optimizes for its specified objective. The Human layer ensures that the objective remains aligned with human welfare. The Root layer establishes what human welfare means in the given context. This three-layer structure is not a feature that can be added to an existing system; it is the organizing principle of the system itself.
